\documentclass[12pt]{article}
\usepackage[utf8]{inputenc}
\usepackage[T1]{fontenc}
\usepackage{hyperref}
\title{Southern Toussian Language Resources for K-12 Curriculum Development}
\author{Compiled by Claude Sonnet 5.5 (claude-sonnet-5-5) for Andr\'as Kornai}
\date{}
\begin{document}
\maketitle

\section{Introduction}
Southern Toussian (ISO 639-3 wib; also Win, Tusia, Tusian) is spoken in and around Toussiana in southwestern Burkina Faso. Wikipedia treats it as the southern dialect of Tusya, with Northern Toussian carrying the code tsp. Speaker counts differ by source: about 19,500 (a 2007 SIL figure), about 39,000 (1995), and 46,000 (Joshua Project). Its place within Niger-Congo is uncertain. Wikitongues gives an EGIDS rating of 6b (threatened). Search results showed few educational resources.

\section{Primary Educational Resources}
No primers, textbooks or school programmes were found in search results.

\section{Linguistic and Reference Materials}
The Wikipedia article on Tusya gives an overview of the two dialects: \url{https://en.wikipedia.org/wiki/Tusya_language}.
The Max Planck number-system database has a page with data for Southern Toussian: \url{https://lingweb.eva.mpg.de/channumerals/Toussian-Southern.htm}.
WALS has an entry for the language: \url{https://wals.info/languoid/lect/wals_code_tou}.
No dictionary or grammar was found.

\section{Religious and Community Resources}
A New Testament in Southern Toussian was published in 2013 (copyright Wycliffe Bible Translators, in cooperation with ANTBA). It is available at eBible.org: \url{https://ebible.org/study/content/texts/wib/about.html}, with download details at \url{https://ebible.org/find/show.php?id=wib}. YouVersion carries it as ``Nouveau Testament en toussian'': \url{https://bible.com/versions/1311}.
Scripture Earth also lists it: \url{https://www.scriptureearth.org/00eng.php?iso=wib}.
Joshua Project profile: \url{https://joshuaproject.net/people_groups/19274/UV}.

\section{Digital and Audio Resources}
The New Testament is offered as mobile HTML, ePub, Kindle, PDF and a Sword module through eBible.org. Wikitongues has a language page: \url{https://wikitongues.org/languages/wib/}. No audio or video was found.

\section{Curriculum Development Context}
The language is still used by its community, but little educational material exists. The New Testament is the main written text and could be a source for orthography and reading texts. A school curriculum would need an orthography review, a lexicon, and locally written readers, with local educators and the national literacy bodies. French is the usual school language, so a bilingual approach would be natural.
\end{document}